\documentclass[13pt]{article}

\usepackage{amsmath}
\usepackage{mathtools}
\usepackage[utf8]{inputenc}
\usepackage[english,greek]{babel}
\usepackage{graphicx}
\graphicspath{{media/}}
\usepackage[
  a4paper,
  top=2.5cm,
  bottom=2.5cm,
  left=2.5cm,
  right=2.5cm
]{geometry}
\usepackage{siunitx}
\usepackage{xcolor}
\usepackage{pdfpages}
\usepackage{soul}
\usepackage{listings}
\usepackage{subcaption}
\usepackage{caption}

\lstset{%
breaklines=true,
breakatwhitespace=true, 
%columns=fullflexible
}

\newcommand{\eng}[1]{\foreignlanguage{english}{#1}}
\newcommand{\langeng}{\selectlanguage{english}}
\newcommand{\langgrc}{\selectlanguage{greek}}

\title{Λειτουργικά Συστήματα Υπολογιστών\\[1.2ex]
        \Large Γρίφος}
\date{\today}
\author{Ίωνας Τσακίρης 03123077\footnote{\eng{el23077@mail.ntua.gr}}}


% Definition of \maketitle
\makeatletter         
\def\@maketitle{
\raggedright
\mbox{}\hfill
\begin{minipage}{2cm}
    \includegraphics[width = 20mm,valign=b]{pyrforos}
\end{minipage}
\hfill
\begin{minipage}{12cm}
    {\Large \bfseries Εθνικό Μετσόβιο Πολυτεχνείο}\\ 
    {\Large Σχολή Ηλεκτρολόγων Μηχανικών και Μηχανικών\\Υπολογιστών}
\end{minipage}

\\[16ex]
\begin{center}
{\Huge \bfseries \@title }\\[4ex] 
{\Large  \@author}\\[4ex] 
\@date\\[8ex]
\end{center}}
\makeatother

\begin{document}
\pagenumbering{gobble}
\maketitle
\newpage
\pagenumbering{arabic}
\section*{\ul{Εμπόδιο 0}}
\subsection*{\ul{Διάγνωση}}
Χρησιμοποιόντας την \eng{\texttt{ltrace -Sf}} βλέπουμε την κλήση \eng{\texttt{[pid 10537] open(".hello_there", 0, 00)             = -1}}. 

\subsection*{\ul{Επίλυση}}
Δημιουργούμε το ζητούμενο αρχείο ώστε η κλήση να επιτυγχάνει.

\vspace{5mm}

\section*{\ul{Εμπόδιο 1}}
\subsection*{\ul{Διάγνωση}}
Όμοια, γίνεται κλήση της \eng{open(..., O_WRONLY, ...)}.


\subsection*{\ul{Επίλυση}}
Αλλάζουμε την επίτρεψη του αρχείου σε \texttt{0444}, δηλαδή σε \eng{read-only}.
\vspace{5mm}

\section*{\ul{Εμπόδιο 1}}
\subsection*{\ul{Διάγνωση}}
\subsection*{\ul{Επίλυση}}

\langeng
\begin{lstlisting}[language=c]
\end{lstlisting}
\langgrc



\end{document}
